\section{Distributions de Bose--Einstein et de Fermi--Dirac}
\label{v:rec:ocupacion:section:02}

La formulation d’équilibre part du Hamiltonien et du nombre total
d’occupations,
\begin{equation}
 H=\sum_kE_k^{\rm occ} n_k,
 \qquad
 N=\sum_k n_k,
 \qquad
 K_\mu:=H-\mu N.
 \label{v:rec:eq:hamiltoniano-gran-canonico-rev3}
\end{equation}
Pour \(\beta>0\), l’état grand canonique est
\begin{equation}
 \rho_{\beta,\mu}
 =\frac{e^{-\beta K_\mu}}{\mathcal Z_{\beta,\mu}},
 \qquad
 \mathcal Z_{\beta,\mu}
 =\operatorname{Tr}(e^{-\beta K_\mu}).
 \label{v:rec:eq:estado-gran-canonico-rev3}
\end{equation}

\begin{theorem}[Factorisation grand canonique et lois d’occupation]
\label{v:rec:thm:factorizacion-gran-canonica-rev3}
Pour des niveaux d’énergie \(E_k^{\rm occ}\) de dégénérescences \(g_k\), la
trace fermionique se factorise comme
\begin{equation}
 \mathcal Z_{\rm FD}
 =\prod_k\bigl(1+e^{-\beta(E_k^{\rm occ}-\mu)}\bigr)^{g_k},
 \qquad
 \langle n_k\rangle_{\rm FD}
 =\frac{g_k}{e^{\beta(E_k^{\rm occ}-\mu)}+1}.
 \label{v:rec:eq:factorizacion-fd-rev3}
\end{equation}
Dans le secteur bosonique, sous la condition
\(\mu<\inf_kE_k^{\rm occ}\),
\begin{equation}
 \mathcal Z_{\rm BE}
 =\prod_k\bigl(1-e^{-\beta(E_k^{\rm occ}-\mu)}\bigr)^{-g_k},
 \qquad
 \langle n_k\rangle_{\rm BE}
 =\frac{g_k}{e^{\beta(E_k^{\rm occ}-\mu)}-1}.
 \label{v:rec:eq:factorizacion-be-rev3}
\end{equation}
\end{theorem}

\begin{proof}
L’opérateur \(K_\mu\) est une somme d’opérateurs d’occupation commutants, de
sorte que sa trace se factorise par modes. Dans chaque mode fermionique, on somme les
occupations \(0\) et \(1\), et le facteur correspondant est
\(1+u_k\), où
\(u_k=e^{-\beta(E_k^{\rm occ}-\mu)}\). Dans chaque mode bosonique, on somme toutes
les occupations non négatives et on obtient la série géométrique
\((1-u_k)^{-1}\) ; la condition sur \(\mu\) garantit \(u_k<1\). Les
puissances \(g_k\) incorporent la dégénérescence. Enfin,
\[
 \langle n_k\rangle
 =-\frac1\beta\frac{\partial}{\partial E_k^{\rm occ}}
   \log\mathcal Z
\]
produit les deux lois d’occupation
\cite{Bose1924,Einstein1924,Fermi1926,Dirac1926,Pathria2011}.
\end{proof}
